\documentclass[12pt]{article}
\usepackage[margin=1in]{geometry}

\title{README\\FlightSound}
\author{Group 05}
\date{May 10, 2026}

\begin{document}

\maketitle
\newpage

\tableofcontents
\newpage

\section*{Version History}
\addcontentsline{toc}{section}{Version History}

\begin{tabular}{|l|l|l|l|}
\hline
Version & Date & Author & Description \\
\hline
1.0 & May 10, 2026 & Group 06 & Initial draft \\
\hline
\end{tabular}

\newpage

\section{Project Management}

The project board for FlightSound is managed through Jira. All sprint tasks, bug reports, and team assignments can be found at the following link:

https://calstatela-group5-final.atlassian.net/jira/software/projects/SCRUM/boards/1?atlOrigin=eyJpIjoiNDMxNWJlMGNjOTM1NDAwYmI0MmQ1NjQzMGU4OWE5MzYiLCJwIjoiaiJ9

\newpage

\section{Objective}

FlightSound is a web application built as a course project for CS3338. The goal of the project is to display real-time flight activity on an interactive map and estimate how loud nearby aircraft are based on their altitude and aircraft type. The backend is written in Java and connects to a MySQL database. The frontend is built with HTML, CSS, and JavaScript using the Google Maps API to render the map.

The core objectives of the project are:

\begin{itemize}
\item Retrieve live flight data from the FlightLabs API and display aircraft positions on an interactive Google Maps interface.
\item Estimate noise levels for each aircraft using a simplified formula based on altitude and aircraft type.
\item Display current weather conditions for the selected airport using the Open-Meteo API.
\item Store flight, weather, and noise data in a MySQL database.
\item Provide a simple and easy to navigate interface for users to find and interact with flight data.
\end{itemize}

\newpage

\section{Goals and Purpose}

\subsection{Why FlightSound Matters}
People who live near airports often have no easy way to know when aircraft noise will be at its worst. FlightSound gives users a simple way to see which flights are nearby, how loud they are estimated to be, and what the current weather conditions are at a selected airport. The idea is that someone could open the application, select their nearest airport, and use the information to plan around noisy periods in their day.

\subsection{Goals}
\begin{itemize}
\item Build a working full-stack web application using Java, MySQL, and external REST APIs
\item Give the team practical experience connecting a frontend, backend, database, and multiple third-party APIs into one functioning system.
\item Produce clean and maintainable code backed by proper documentation including this README, the SRS, and the SDD.
\item Deliver a project that could be extended with additional features in the future such as historical noise data, user accounts, or flight path predictions.
\end{itemize}

\newpage

\section{How to Access and Download}

\subsection{Accessing the Application}
FlightSound can be accessed through any modern web browser. Navigate to the hosted URL provided by the deployment environment. No installation is required to use the application.

\subsection{Downloading and Running Locally}
To run FlightSound locally, follow the steps below:

\begin{enumerate}
\item Make sure Java 17 or higher and MySQL 8.x are installed on your machine.
\item Clone the project repository from GitHub.
\item Navigate to the project directory in your terminal.
\item Set up a configuration file with the required environment variables: your FlightLabs API key, Google Maps API key, and MySQL database credentials.
\item Run the database setup to create the required tables: Airport, Flight, Weather, NoiseImpact, and Query.
\item Compile and start the Java backend server.
\item Open the frontend files from the frontend directory in your browser or serve them through a local web server.
\item The application will be available at the local address provided by the server.
\end{enumerate}

\end{document}
